\documentclass[11pt,a4paper]{article}

\usepackage[margin=1in]{geometry}
\usepackage{amsmath,amssymb}
\usepackage{graphicx}
\usepackage{booktabs}
\usepackage{array}
\usepackage[colorlinks=true,linkcolor=blue,citecolor=blue,urlcolor=blue]{hyperref}
\usepackage{caption}
\usepackage{enumitem}
\usepackage{xcolor}

% ---------------------------------------------------------------
%  Placeholder styling for the forthcoming hardware section
% ---------------------------------------------------------------
\newcommand{\hwplaceholder}[1]{%
  \begin{center}
    \colorbox{gray!15}{\parbox{0.92\textwidth}{%
      \vspace{0.5em}
      {\color{gray!60!black}\textbf{[PLACEHOLDER --- to be completed with real-hardware results]}}\\[4pt]
      {\color{gray!55!black}#1}
      \vspace{0.4em}
    }}
  \end{center}
}

\title{Quantum Machine Learning for Pneumonia Detection in\\
Chest X-Ray Imaging: an Illustrative Case Study}

\author{Author names\\[2pt] \small Affiliation to be inserted}

\date{}

\begin{document}

\maketitle

% ==============================================================
%  ABSTRACT
% ==============================================================
\begin{abstract}
Deep convolutional neural networks (CNNs) have become the de facto standard
for automated pneumonia screening from chest X-ray (CXR) images, routinely
reporting diagnostic accuracies above 85\%.  This performance, however, is
obtained at the cost of classification heads with hundreds of thousands of
trainable parameters, motivating interest in parameter-efficient alternatives.
Variational quantum circuits (VQCs), operating within the constraints of Noisy
Intermediate-Scale Quantum (NISQ) hardware, have been proposed as compact
decision layers that can be appended to classical feature extractors.  Beyond
the pure VQC paradigm, quantum convolutional neural networks (QCNNs) and
quantum kernel methods have been explored as complementary architectural
routes.  This paper reviews the computational and algorithmic foundations of
such hybrid quantum--classical architectures for CXR-based pneumonia detection,
covering classical CNN baselines, quantum data-encoding strategies (angle,
amplitude, IQP feature maps, and data re-uploading), variational and
convolutional ansatz design, quantum kernel construction, and the
barren-plateau phenomenon that limits the trainability of deep parametrized
circuits together with its principal mitigation strategies.  As a concrete
illustration, we describe and analyze a hybrid pipeline that couples a frozen
ConvNeXt-Tiny feature extractor with nonlinear autoencoder dimensionality
reduction (768$\rightarrow$64) and a 6-qubit, 54-parameter data re-uploading
VQC classifier, evaluated on a publicly available pediatric CXR dataset under
an explicit train/test class-prevalence shift.  The hybrid model attains a
test AUC-ROC of 0.860 (versus 0.940 for a classical multilayer perceptron
trained on identical features) while using approximately $39\times$ fewer
trainable parameters in its decision layer (54 versus 2{,}113).  We situate
this case study within a broader survey of
related hybrid approaches for respiratory and cardiac CXR analysis reported
in the literature, and ground a discussion of the algorithmic trade-offs,
evaluation methodology, and open problems---including state-preparation
overhead, barren plateaus, dataset shift, and the absence, at present circuit
scales, of a demonstrated computational quantum advantage.  A dedicated
section is reserved for forthcoming experiments on physical quantum hardware
with error mitigation; its completion will allow direct quantification of the
simulator-to-hardware performance gap.
\end{abstract}

\noindent\textbf{Keywords:} quantum machine learning; variational quantum
circuits; hybrid quantum--classical neural networks; quantum convolutional
neural networks; quantum kernel methods; pneumonia detection; chest X-ray
classification; NISQ; barren plateaus; medical image analysis

% ==============================================================
\section{Introduction}
\label{sec:intro}
% ==============================================================

Pneumonia remains one of the leading causes of morbidity and mortality
worldwide, with an estimated global burden in the range of millions of cases
annually and disproportionate impact on children under five and
immunocompromised adults \cite{cdc2025,iqwig2021}.  Chest radiography (CXR)
is the most widely used first-line imaging modality for pneumonia screening
because it is inexpensive, fast, and broadly available compared to computed
tomography \cite{iqwig2021}.  Manual interpretation of CXR images is subject
to inter-observer variability, reader fatigue, and uneven access to trained
radiologists, which has motivated three decades of research into
computer-aided diagnosis and, more recently, deep-learning-based automated
screening \cite{naz2020,kundu2021}.

CNNs, and in particular deep residual architectures such as ResNet, have
demonstrated diagnostic performance on CXR pneumonia classification that is
competitive with expert radiologists, typically exceeding 85\% classification
accuracy \cite{he2016,kundu2021,ortiz2022}.  This performance,
however, comes at a substantial computational cost: fully connected
classification heads appended to CNN backbones commonly require on the order
of $10^5$--$10^6$ trainable parameters.  The associated memory and energy
costs are non-trivial for edge or point-of-care deployment scenarios and
motivate the search for parameter-efficient decision layers.

In parallel, the maturation of NISQ devices has generated interest in
variational quantum algorithms (VQAs) as parameter-efficient function
approximators \cite{schuld2018,preskill2018}.  A VQC parametrizes a family of
unitary transformations $U(\theta)$ and, when combined with classical
optimization of $\theta$, can implement expressive decision boundaries using
a number of trainable parameters that grows only linearly with circuit depth
and qubit count \cite{mitarai2018,biamonte2017}.  This has prompted a growing
body of work on \emph{hybrid quantum--classical} (HQC) architectures, in which
a classical CNN performs high-dimensional feature extraction and a quantum
component performs low-dimensional classification, confining the quantum
computation to regimes tractable on current simulators or hardware.

A number of recent studies have applied HQC architectures specifically to
thoracic imaging tasks.  Kulkarni et al.~\cite{kulkarni2023} demonstrated that
replacing a classical CNN layer with a VQC yields statistically significant
improvements for pneumonia detection on CXR images.  Rao et
al.~\cite{rao2024} reported a 98.1\% test accuracy on a COVID-19 Radiography
Dataset using a hybrid CNN--VQC pipeline deployed partly on an IBM~Q simulator,
while Decoodt et al.~\cite{decoodt2023} achieved an AUC of up to 0.93 for
cardiomegaly detection using a DenseNet-121 backbone augmented with a
variational quantum layer.  At the same time, attention has shifted beyond the
``quantum classifier head'' paradigm: quantum convolutional neural networks
(QCNNs) designed to mimic spatial locality in quantum circuits
\cite{cong2019,wei2022}, and quantum kernel methods that compute high-dimensional
inner products via quantum interference \cite{havlicek2019,haug2021}, represent
architecturally distinct alternatives that may better exploit the inductive
biases of image data.

Despite this activity, empirical evidence quantifying the accuracy--efficiency
trade-off of HQC pipelines on realistic medical imaging tasks---particularly
under distribution shift between training and deployment data---remains
limited.  Most published studies report results on a single static
train/test split, on noise-free simulators, and without explicit discussion of
the state-preparation and trainability bottlenecks intrinsic to variational
quantum computing \cite{sharma2023survey}.

\subsection{Contributions and Scope}

This paper has two purposes.  First, it provides a structured,
computationally oriented review of the building blocks required to construct
hybrid quantum--classical pneumonia classifiers: classical CNN backbones,
dimensionality reduction, quantum data-encoding strategies, variational and
convolutional ansatz design, quantum kernel construction, and the
barren-plateau problem including its principal mitigation strategies.  Second,
it presents a detailed case study---a ConvNeXt-Tiny/autoencoder/VQC pipeline
evaluated on a public pediatric CXR dataset---that grounds the algorithmic
discussion
with explicit numbers and illustrates the current gap between classical and
quantum-assisted approaches.

The clinical and radiological background is kept intentionally brief; the
emphasis throughout is on encoding schemes, circuit design, trainability,
computational complexity, and evaluation methodology, consistent with a
computer-science- and quantum-technology-oriented readership.  We deliberately
avoid overstating the clinical readiness of the reviewed methods: at the qubit
counts considered here, no computational quantum advantage is established, and
the results should be interpreted as a feasibility and parameter-efficiency
study rather than as a validated diagnostic tool.  A separate section
(Section~\ref{sec:hardware}) is reserved for forthcoming experimental results
on physical quantum hardware; its inclusion reflects the acknowledged
importance of closing the simulator-to-hardware gap and will complete the
empirical contribution of this work.

The remainder of the paper is organized as follows.
Section~\ref{sec:background} introduces the clinical context and the classical
and quantum computational preliminaries.  Section~\ref{sec:classical} reviews
classical CNN-based approaches to CXR pneumonia detection.
Section~\ref{sec:qml} describes quantum data encoding, variational circuits,
QCNNs, quantum kernels, and the barren-plateau phenomenon.
Section~\ref{sec:hybrid} surveys related HQC work for CXR analysis and
presents the ConvNeXt-Tiny/autoencoder/VQC case study.
Section~\ref{sec:hardware}
provides the placeholder for real-hardware experiments.
Section~\ref{sec:eval} discusses experimental design and evaluation metrics.
Section~\ref{sec:limitations} enumerates open problems and limitations.
Section~\ref{sec:future} outlines future research directions, and
Section~\ref{sec:conclusion} concludes.

% ==============================================================
\section{Background}
\label{sec:background}
% ==============================================================

\subsection{Clinical and Imaging Context}

Pneumonia is an acute inflammatory condition of the pulmonary parenchyma,
most commonly of bacterial or viral etiology, in which the alveoli fill with
fluid or exudate, impairing gas exchange \cite{iqwig2021}.  On a CXR, healthy
air-filled lung tissue appears radiolucent (dark), whereas pneumonic
consolidation, interstitial infiltrates, and air bronchograms appear as
localized or diffuse radiopaque (light) regions \cite{iqwig2021}.  Bacterial
pneumonia (e.g., \emph{Streptococcus pneumoniae}) typically presents with
lobar consolidation, while viral pneumonia tends to produce more diffuse
interstitial patterns \cite{iqwig2021}.  This visual heterogeneity, combined
with overlapping appearances with other conditions (pulmonary edema,
atelectasis, fibrosis), contributes to inter-observer variability in manual
reading and motivates automated decision support \cite{naz2020}.

\subsection{Classical Deep Learning for Medical Image Classification}

Modern CXR classifiers are almost universally built on deep CNN backbones
pretrained on large natural-image corpora (e.g., ImageNet) and subsequently
fine-tuned or used as fixed feature extractors on medical imaging tasks.
ResNet architectures address the vanishing-gradient and degradation problems
that historically limited the trainability of very deep networks by introducing
residual (skip) connections: instead of learning a direct mapping $H(x)$, each
residual block learns a residual function $F(x)$ such that
\begin{equation}
  H(x) = F(x) + x,
  \label{eq:resnet}
\end{equation}
which allows gradients to propagate across many layers without attenuation
\cite{he2016}.  ResNet-50 uses a bottleneck design in which each residual block
consists of a $1\times1$ convolution (dimensionality reduction), a $3\times3$
convolution (spatial filtering), and a further $1\times1$ convolution
(dimensionality restoration), substantially reducing parameter count relative
to a naive stacking of $3\times3$ convolutions while preserving
representational depth \cite{he2016}.  In a typical CXR pipeline, the network
is used \emph{headless}, i.e., without its final classification layer,
producing a fixed-length embedding (commonly of dimension 2048 for ResNet-50)
for each input image, which is then passed to a task-specific classifier.

\subsection{Elements of Quantum Computation}

A qubit is the basic unit of quantum information, represented as a normalized
vector in a two-dimensional complex Hilbert space,
\begin{equation}
  \lvert \psi \rangle = \alpha \lvert 0 \rangle + \beta \lvert 1 \rangle,
  \qquad \lvert \alpha \rvert^2 + \lvert \beta \rvert^2 = 1,
  \label{eq:qubit}
\end{equation}
where $\alpha,\beta \in \mathbb{C}$ are probability amplitudes and
$\{\lvert 0\rangle, \lvert 1\rangle\}$ is the computational basis
\cite{nielsen2010}.  Geometrically, a pure single-qubit state can be
represented as a point on the Bloch sphere,
\begin{equation}
  \lvert \psi \rangle = \cos\!\left(\frac{\theta}{2}\right)\lvert 0 \rangle
  + e^{i\varphi}\sin\!\left(\frac{\theta}{2}\right)\lvert 1 \rangle,
  \label{eq:bloch}
\end{equation}
and training a variational circuit can be understood, at the single-qubit
level, as searching for rotation angles that steer this state toward a
measurement outcome associated with the desired class.

Quantum algorithms are composed of unitary \emph{gates}.  Single-qubit
rotation gates $R_x, R_y, R_z$ (or the three-angle rotation
$\mathrm{Rot}(\phi,\theta,\omega)$) rotate the state around the
corresponding axis of the Bloch sphere and constitute the trainable degrees of
freedom of a variational circuit; the Hadamard gate creates equal
superpositions; and the two-qubit controlled-NOT (CNOT) gate is the standard
primitive for generating \emph{entanglement}, a distinctively quantum
correlation between qubits that cannot be reproduced by any local, classical
description of the joint system and that is conjectured to underlie the
representational power of quantum models \cite{nielsen2010}.

The \emph{expressibility} of a parametrized circuit is a quantitative measure
of how uniformly it can generate states over the Hilbert space, often defined
via the deviation of the induced distribution from the Haar measure
\cite{sim2019}.  Expressibility and entanglement capacity jointly govern the
circuit's ability to represent complex functions and, as detailed in
Section~\ref{sec:barren}, its susceptibility to trainability issues.

\subsection{The NISQ Era and Variational Quantum Algorithms}

Contemporary quantum processors operate in what Preskill termed the Noisy
Intermediate-Scale Quantum (NISQ) era: devices with tens to a few hundred
physical qubits that are not yet fault-tolerant and remain subject to gate
errors and decoherence \cite{preskill2018}.  Variational quantum algorithms
are the dominant algorithmic paradigm for extracting practical value from NISQ
hardware.  A VQA defines a parametrized unitary $U(\theta)$, executes it on
a (possibly classically simulated) quantum device, measures an observable, and
updates $\theta$ using a classical optimizer---a hybrid quantum--classical
optimization loop that keeps quantum resource requirements shallow while
delegating gradient-based optimization to classical compute
\cite{schuld2018,mitarai2018}.  Variational Quantum Classifiers (VQCs) are the
supervised-learning instance of this paradigm and form the quantum component
of the hybrid architectures discussed in this paper.

% ==============================================================
\section{Classical Approaches to Pneumonia Detection from Chest X-Rays}
\label{sec:classical}
% ==============================================================

The dominant classical paradigm for CXR-based pneumonia detection follows a
transfer-learning recipe: (i)~a CNN backbone pretrained on a large
natural-image dataset is used, either fine-tuned end-to-end or frozen as a
fixed feature extractor; (ii)~a fully connected classification head, typically
consisting of one or more dense layers followed by a sigmoid or softmax
output, is trained on the target task; and (iii)~standard image preprocessing
(resizing, normalization with respect to the pretraining dataset's channel
statistics) is applied to bridge the domain gap between natural and medical
images.  ResNet, DenseNet, and, more recently, vision-transformer backbones
have all been applied to the publicly available pediatric CXR pneumonia dataset
originating from the Guangzhou Women and Children's Medical Center, with
reported accuracies generally exceeding 85\% and F1 scores in the
0.85--0.92 range \cite{kundu2021,ortiz2022,naz2020}.

From a purely computational standpoint, the classification head is the
component of the pipeline most directly comparable to a quantum decision
layer.  For a 64-dimensional input (e.g., after dimensionality reduction,
see Section~\ref{sec:hybrid}) and a single hidden layer of width 32, a
standard multilayer perceptron (MLP) classification head requires
\begin{equation}
  P_{\text{MLP}} = (64 \times 32 + 32) + (32 \times 1 + 1) = 2{,}080 + 33
  = 2{,}113 \text{ parameters},
  \label{eq:mlp_params}
\end{equation}
illustrating that even a comparatively shallow classifier on a compressed
feature space carries a non-trivial parameter budget.  This observation is the
principal motivation for exploring VQC-based decision layers, whose parameter
count scales as $O(nL)$ in the number of qubits $n$ and circuit layers $L$,
rather than with the width of a dense layer (Section~\ref{sec:qml}).

Two further classical design choices are directly relevant to the hybrid
architectures discussed below.  First, class imbalance---pneumonia typically
constitutes 70--75\% of samples in public pediatric CXR datasets---motivates
the use of weighted loss functions (e.g., weighted binary cross-entropy) rather
than unweighted cross-entropy, to avoid a majority-class bias.  Second,
dimensionality reduction (via PCA, linear discriminant analysis, or univariate
feature selection) is frequently required before any quantum processing step,
since CNN embeddings (commonly 512--2048-dimensional) exceed the practical
qubit budgets of NISQ-era hardware and simulators by orders of magnitude.

% ==============================================================
\section{Quantum Machine Learning Approaches}
\label{sec:qml}
% ==============================================================

\subsection{Data Encoding Strategies}
\label{sec:encoding}

Before a classical feature vector $x \in \mathbb{R}^N$ can be processed by a
quantum circuit, it must be mapped, via a \emph{state-preparation} routine,
onto a quantum state $\lvert \psi(x) \rangle$ in an $n$-qubit Hilbert space.
The choice of encoding is consequential: it determines the geometry of the
induced kernel, the depth of the state-preparation circuit, and, as recent
theoretical work has shown, the trainability of the subsequent variational
circuit \cite{kashif2023,schuld2021kernel}.

\paragraph{Basis encoding} maps a bit string directly onto computational basis
states, requiring one qubit per bit of the (typically pre-discretized) classical
input.  It scales poorly for continuous, high-dimensional features and is
rarely used for image-derived embeddings.

\paragraph{Angle (rotation) encoding} assigns each classical feature $x_i$ to
a rotation angle on a dedicated qubit, e.g., $R_y(x_i)$, and therefore
requires $n = N$ qubits for an $N$-dimensional feature vector.  It has the
practical advantage of a shallow, easily implementable preparation circuit and
has been shown to produce more consistent results on real QPUs than more
complex encodings in at least one comparative hardware study \cite{qkernel2022}.

\paragraph{Amplitude encoding} embeds a normalized vector
$x \in \mathbb{R}^N$, $\lVert x \rVert_2 = 1$, into the amplitudes of an
$n$-qubit state,
\begin{equation}
  \lvert \psi(x) \rangle = \sum_{i=0}^{N-1} x_i \lvert i \rangle
  \in \mathcal{H}_{2^n}, \qquad N = 2^n,
  \label{eq:amplitude}
\end{equation}
requiring only $n = \log_2 N$ qubits.  This exponential compression is the
primary motivation for its use in hybrid CNN--VQC pipelines, since it allows,
e.g., a 64-dimensional PCA-reduced embedding to be encoded on just
$\log_2 64 = 6$ qubits.

\paragraph{IQP and ZZ feature maps} apply layers of classically controlled
phase rotations---in particular the instantaneous quantum polynomial (IQP)
circuit structure---to produce highly non-linear, problem-specific feature maps
\cite{havlicek2019}.  The resulting kernel $k(x, x') =
|\langle \psi(x) | \psi(x') \rangle|^2$ can be exponentially hard to
evaluate classically, lending potential computational interest to quantum
kernel classifiers (Section~\ref{sec:kernels}).  Empirically, IQP-based
encodings have been observed to improve precision on some tasks while reducing
recall and $F_1$ scores relative to angle encoding, and their performance is
more sensitive to hardware noise \cite{qkernel2022}.

\paragraph{Data re-uploading} addresses the limited expressibility of a
single encoding layer by repeatedly interleaving data-encoding gates with
trainable rotation layers throughout the circuit, effectively allowing each
qubit to ``see'' the input multiple times \cite{perez2020}.  This strategy
has been shown to confer universal function-approximation properties to single
qubits and provides a practical route to expressible, shallow circuits
compatible with NISQ depth budgets.

The compression advantage of amplitude encoding is not free: in the general
case, exact amplitude-encoding state preparation for an arbitrary
$N$-dimensional vector requires a circuit of depth $O(N)$, i.e.,
\emph{exponential} in the number of qubits, unless the vector possesses
exploitable structure or approximate preparation is acceptable
\cite{aaronson2015}.  This ``input bottleneck'' is a widely recognized caveat
in the QML literature, and claims of exponential quantum speedup based on
amplitude encoding alone should be treated with caution
\cite{aaronson2015,biamonte2017}.  In practice, most simulation-based
studies---including the case study reviewed in Section~\ref{sec:hybrid}---rely
on software-level amplitude-embedding routines (e.g., PennyLane's
\texttt{AmplitudeEmbedding}) whose classical preprocessing cost is not charged
against the ``quantum'' resource budget, a simplification that must be kept
in mind when extrapolating simulator-based results to real hardware.

\subsection{Variational Quantum Circuits}
\label{sec:vqc}

Given an encoded input state $\lvert \psi(x) \rangle$, a VQC applies a
parametrized unitary $U(\theta)$---the \emph{ansatz}---and measures an
observable $M$ (typically a Pauli operator) to produce a scalar prediction,
\begin{equation}
  f(x, \theta) = \langle \psi(x) \rvert\, U^{\dagger}(\theta)\, M\,
  U(\theta) \lvert \psi(x) \rangle \in [-1, 1].
  \label{eq:vqc}
\end{equation}
When $M = Z_0$ (the Pauli-$Z$ operator on the first qubit), this quantity is
linearly remapped to a pseudo-probability,
\begin{equation}
  \sigma(f) = \frac{f + 1}{2} \in [0, 1],
  \label{eq:sigmoid_remap}
\end{equation}
which can be thresholded for binary classification.  A common and
hardware-efficient ansatz design alternates layers of single-qubit rotations
with a fixed entangling pattern: each layer $\ell = 1,\dots,L$ applies a
general single-qubit rotation $\mathrm{Rot}(\phi,\theta,\omega)$ to every
qubit, followed by a ring of CNOT gates connecting qubit $w$ to qubit
$(w+1) \bmod n$.  For $n$ qubits and $L$ layers, this yields
\begin{equation}
  P_{\text{VQC}} = 3\, n\, L
  \label{eq:vqc_params}
\end{equation}
trainable parameters---e.g., $3 \times 6 \times 3 = 54$ parameters for a
6-qubit, 3-layer circuit, versus the $2{,}113$ parameters of a
comparable classical dense head (Equation~\ref{eq:mlp_params}).  Gradients
of $f(x,\theta)$ with respect to $\theta$ can be computed on real hardware
via the parameter-shift rule, or, in simulation, via reverse-mode automatic
differentiation \cite{mitarai2018}.

\subsection{Quantum Convolutional Neural Networks}
\label{sec:qcnn}

Quantum convolutional neural networks (QCNNs) translate the locality and
parameter sharing of classical CNNs into the quantum circuit formalism.  The
archetypal QCNN design, introduced by Cong et al.~\cite{cong2019} and adapted
for NISQ hardware by Wei et al.~\cite{wei2022}, alternates \emph{quantum
convolutional layers}---parametrized two-qubit gates applied to neighboring
qubit pairs---with \emph{quantum pooling layers} that halve the active qubit
register via measurements and classically controlled single-qubit operations,
mimicking the spatial downsampling of classical max-pooling.  A key theoretical
result is that QCNNs built on translational symmetry are free from the barren-plateau problem
under local cost functions, providing a trainability guarantee absent from
more generic parametrized circuits \cite{pesah2021}.

In the context of medical image classification, Khatoniar et
al.~\cite{khatoniar2024} applied QCNN variants to CXR pneumonia classification
under extreme feature reduction, reporting competitive robustness relative to
VQC baselines of comparable depth.  Wei et al.~\cite{wei2022} demonstrated
digit recognition and image filtering tasks using a hardware-implementable QCNN
whose parameter count is \emph{independent of input size} (determined only by
circuit depth and qubit count), an attractive property for deployment on
resource-constrained platforms.

It is important to distinguish QCNNs from \emph{quanvolutional} layers
\cite{henderson2020quanvolutional}, in which a small fixed or parametrized
quantum circuit processes local image patches in a sliding-window fashion,
analogously to a classical convolutional kernel, before passing the resulting
feature map to a standard classical CNN.  Quanvolutional layers add a quantum
pre-processing stage rather than replacing the classical network, and are
primarily motivated by the conjecture that quantum circuits can generate
feature transformations qualitatively distinct from those achievable by
classical filters.

\subsection{Quantum Kernel Methods}
\label{sec:kernels}

Quantum kernel methods provide an alternative to the parametrized-circuit
paradigm by using a quantum circuit exclusively to evaluate a kernel function,
\begin{equation}
  k(x, x') = \lvert \langle \psi(x) \lvert \psi(x') \rangle \rvert^2,
  \label{eq:kernel}
\end{equation}
and delegating classification to a classical support vector machine (SVM) or
kernel ridge regressor trained on the resulting kernel matrix.  Havl\'i\v{c}ek
et al.~\cite{havlicek2019} showed, using the ZZ feature map, that the
corresponding quantum kernel can be classically hard to estimate exactly,
providing a theoretical grounding for potential quantum advantage.  Schuld and
Killoran~\cite{schuld2019kernel} established the connection between quantum
models and kernel methods more broadly, showing that every VQC implicitly
defines a kernel via its state overlap.

For medical imaging, the high dimensionality of image-derived feature vectors
poses a practical challenge: evaluating the kernel matrix over a training set
of $N_{\text{tr}}$ samples requires $O(N_{\text{tr}}^2)$ circuit executions,
each querying the state overlap of two encoded feature vectors.  Haug et
al.~\cite{haug2021} proposed scalable quantum kernel estimation via
randomized measurements, reducing both the circuit depth required and the
sensitivity to noise, and demonstrated large-scale kernel estimation on IBM
quantum hardware with noise mitigation.  In a comparative hardware study,
angle encoding achieved the best overall kernel performance on real QPUs
among the encodings tested, albeit with pronounced run-to-run variability,
while IQP feature maps performed more consistently yet underperformed
relative to classical kernels; dynamical decoupling error mitigation did not
uniformly improve performance and was encoding-dependent
\cite{qkernel2022}, underscoring that hardware-aware encoding selection is
essential rather than optional.

\subsection{The Barren-Plateau Problem}
\label{sec:barren}

The central obstacle to training deep or highly entangling VQCs is the
\emph{barren-plateau} phenomenon: for a broad class of randomly initialized,
sufficiently expressive ansätze evaluated against a global cost function, the
variance of the cost-function gradient with respect to any single parameter
$\theta_i$ decays exponentially with the number of qubits $n$,
\begin{equation}
  \mathrm{Var}\!\left[\partial_{\theta_i} C(\theta)\right]
  = O\!\left(2^{-n}\right),
  \label{eq:barren}
\end{equation}
so that, in expectation, gradient-based optimization becomes exponentially
difficult to distinguish from noise as $n$ grows \cite{mcclean2018}.
Subsequent work has elucidated the mechanisms underlying this concentration.
Cerezo et al.~\cite{cerezo2021} showed that the onset of barren plateaus
depends on the \emph{locality} of the cost function: global cost functions---
those defined via an operator acting non-trivially on all qubits---are markedly
more prone to barren plateaus than local ones.  Ortiz Marrero et
al.~\cite{marrero2021} demonstrated that excessive (volume-law) entanglement in
the circuit is sufficient to induce barren plateaus independently of the
specific ansatz structure.  Kashif and Al-Kuwari~\cite{kashif2023} developed a
unified framework showing that barren-plateau onset depends jointly on data
encoding, ansatz expressibility, and entanglement, providing design constraints
for HQC networks.

These theoretical results have motivated several practical mitigation
strategies:
\begin{itemize}[leftmargin=1.4em]
  \item \textbf{Local cost functions and shallow entangling ansätze.}
    Restricting measurements to local Pauli observables and using
    nearest-neighbor entangling layers (e.g., the CNOT ring in
    Equation~\ref{eq:vqc_params}) substantially suppresses gradient variance
    \cite{cerezo2021}.
  \item \textbf{Controlled initialization.}  Initializing parameters near the
    identity, or near simple product states (``identity blocks''), rather than
    uniformly at random, ensures that the optimization starts in a region of
    non-vanishing gradient \cite{mcclean2018,patti2021}.  Partitioning the
    circuit into cost and non-cost registers and using low-entanglement
    initializations provides a practical recipe \cite{patti2021}.
  \item \textbf{Expressibility-aware circuit design.}  Reducing circuit
    expressibility---via state-efficient ansätze (SEA) \cite{liu2022sea} or
    neural-network-based parameter generation that smooths the loss landscape
    \cite{zhu2024}---trades representational power for trainability.
  \item \textbf{Noise-assisted mitigation.}  Entanglement regularization, the
    addition of controlled Langevin noise, and rotations into the cost
    function's eigenbasis can improve convergence on difficult optimization
    landscapes \cite{patti2021}.
  \item \textbf{Noise mitigation as expressibility reduction.}  Ding et
    al.~\cite{ding2024} showed that adapting circuit expressibility to match
    the noise level of the hardware can act as an error-aware mitigation
    strategy, improving classification accuracy in noisy simulations for
    QCNNs.
\end{itemize}
These design constraints guided the layer-count selection for the case-study
ansatz in Section~\ref{sec:hybrid}: three data re-uploading layers on six
qubits keep the circuit in a regime of moderate expressibility and shallow
entanglement, where gradient-variance concentration is not expected to
dominate training on a noise-free simulator.

\subsection{On the Prospect of Quantum Advantage}
\label{sec:advantage}

A recurring question in the QML literature is whether VQC-based classifiers
can offer a genuine \emph{computational} advantage---in runtime, sample
complexity, or expressivity---over classical models, as opposed to merely a
reduction in the number of trainable parameters.  At the circuit sizes
considered in this review (single-digit to low tens of qubits), no such
advantage has been established for medical image classification: circuits of
this size are efficiently simulable classically, and the reduction in
decision-layer parameter count does not, by itself, imply reduced overall
computational cost, since the classical feature-extraction and
dimensionality-reduction stages typically dominate the total pipeline cost.
Claims of quantum advantage in QML more broadly remain an active and contested
research area, further complicated by the input/output bottleneck discussed
above \cite{aaronson2015,biamonte2017}.  The practical value of hybrid VQC
decision layers at present should therefore be understood primarily in terms
of \emph{parameter efficiency} and \emph{memory footprint}---properties that
may be relevant for edge or resource-constrained deployment once fault-tolerant
or larger-scale NISQ hardware becomes available---rather than as evidence of
an accuracy or speed advantage over well-tuned classical models.

% ==============================================================
\section{Hybrid Quantum--Classical Architectures}
\label{sec:hybrid}
% ==============================================================

\subsection{Design Rationale}

Given the qubit and depth constraints of NISQ hardware
(Section~\ref{sec:qml}) and the high dimensionality of raw medical images, a
natural design principle is to separate \emph{feature learning}, delegated to
a classical CNN operating on the full-resolution image, from \emph{decision
making}, delegated to a compact quantum circuit operating on a low-dimensional,
information-dense representation.  This principle is shared by all hybrid
architectures reviewed below, differing primarily in (i)~the choice of
classical backbone, (ii)~the dimensionality reduction method, (iii)~the
encoding and ansatz design, and (iv)~whether the quantum component serves as a
classifier head, a convolutional feature transformer, or a kernel function.

\subsection{Related Work on Hybrid Quantum--Classical CXR Classifiers}
\label{sec:relwork}

Table~\ref{tab:related} summarizes published hybrid QML studies targeting
chest X-ray analysis.  Three themes emerge from this survey.

\begin{table}[htbp]
\centering
\small
\caption{Representative hybrid quantum--classical studies on CXR or pulmonary
  imaging tasks.  All results are from noise-free simulators unless marked
  with $\dagger$ (IBM~Q backend).  ``HQC-VQC'' denotes a classical feature
  extractor followed by a VQC classifier head; ``HQC-Conv'' denotes a
  quanvolutional or QCNN layer integrated into the CNN pipeline.}
\label{tab:related}
\begin{tabular}{@{}p{3.5cm}p{2.0cm}p{2.0cm}p{1.3cm}p{4.0cm}@{}}
\toprule
Reference & Architecture & Dataset & Task & Key result \\
\midrule
Kulkarni et al.~\cite{kulkarni2023} & HQC-VQC & CXR (Guangzhou) & Pneumonia & Stat.\ significant improvement vs.\ classical CNN baseline \\
Rao et al.~\cite{rao2024} & HQC-VQC ($\dagger$) & COVID-19 Radiography Dataset (15K) & 4-class & 98.1\% test acc.; IBM~Q QASM deployment \\
Decoodt et al.~\cite{decoodt2023} & HQC-VQC & CheXpert (2.4K PA) & Cardiomegaly & AUC $\leq 0.93$; improved Grad-CAM++ localization \\
Sharma~\cite{sharmaR2023hybrid} & HQC-Conv & Custom CXR & Pneumonia & ZFNet + quantum layer hybrid \\
Khatoniar et al.~\cite{khatoniar2024} & QCNN & CXR (Guangzhou) & Pneumonia & Robustness analysis under extreme feature reduction \\
\textbf{This work} & HQC-VQC & CXR (Guangzhou, 5.8K) & Pneumonia & AUC 0.860 (VQC) vs.\ 0.940 (MLP); $39\times$ param.\ reduction \\
\bottomrule
\end{tabular}
\end{table}

\paragraph{Consistent accuracy--efficiency trade-off.}
All surveyed studies confirm the qualitative finding that a hybrid quantum
component can achieve meaningful classification performance at a substantially
reduced decision-layer parameter budget.  However, results that appear
favorable relative to their specific classical baselines often reflect
comparisons against purposely weak baselines rather than state-of-the-art
classical models \cite{aaronson2015}.

\paragraph{Variability in experimental rigor.}
Dataset sizes range from a few hundred to $\sim$15,000 images; train/test
split rationale is frequently underspecified; and class prevalence in the
test set is rarely reported explicitly.  This limits direct cross-study
comparison and motivates the shift-aware evaluation framework presented in
Section~\ref{sec:eval}.

\paragraph{Limited hardware deployment.}
Only Rao et al.~\cite{rao2024} report partial deployment on a real IBM~Q backend
(QASM simulator), and no study provides a systematic comparison of
simulator versus hardware performance under noise.  The gap between
simulator and hardware results---which Section~\ref{sec:hardware} is designed
to quantify---remains the most critical open empirical question in the field.

\subsection{Pipeline Architecture}

The pipeline consists of four stages, summarized in
Table~\ref{tab:pipeline}: (1)~a frozen ConvNeXt-Tiny backbone, used headless,
maps each $224\times224$ input image to a 768-dimensional embedding; (2)~a
domain-adversarial neural network (DANN) with a gradient reversal layer (GRL)
is trained on top of the embedding to suppress the dataset shift between the
training and test distributions (Section~\ref{sec:relwork}); (3)~a
nonlinear autoencoder, fitted on the training embeddings, reduces this
embedding to 64 dimensions, chosen to match the Hilbert-space dimension of a
6-qubit register ($2^6 = 64$); and (4)~a
6-qubit, 3-layer VQC, using amplitude encoding
(Equation~\ref{eq:amplitude}) and the rotation-plus-ring-CNOT ansatz of
Section~\ref{sec:vqc}, produces a scalar pneumonia score via
Equations~\ref{eq:vqc}--\ref{eq:sigmoid_remap}.

\begin{table}[htbp]
\centering
\caption{Hybrid pipeline architecture and parameter counts.}
\label{tab:pipeline}
\begin{tabular}{@{}lllll@{}}
\toprule
Stage & Input & Operation & Output & Parameters \\
\midrule
Feature extraction & $224\times224$ image & ConvNeXt-Tiny (frozen) & 768-D vector & 0 (fixed) \\
Domain adaptation & 768-D & DANN + GRL & 768-D vector & 0 (fixed) \\
Dimensionality reduction & 768-D & Autoencoder (nonlinear) & 64-D vector & 0 (fixed) \\
Quantum classification & 64-D & VQC $U(\theta)$ & score $\in [0,1]$ & 54 \\
\bottomrule
\end{tabular}
\end{table}

Among linear PCA, linear discriminant analysis (LDA), and univariate feature
selection (SelectKBest), a nonlinear autoencoder yielded the most stable
training dynamics and the best validation performance at 64 latent
dimensions; LDA suffered from numerical instability in the
high-dimension/limited-sample regime, and SelectKBest achieved marginally
weaker validation metrics.  The autoencoder therefore replaces the linear
projection used in earlier variants of this pipeline.

\begin{figure}[htbp]
\centering
\fbox{\parbox{0.85\textwidth}{\centering \vspace{1.2em}
\textbf{[Figure placeholder]} \\[4pt]
Chest X-ray $224\times224$ \;$\rightarrow$\; ConvNeXt-Tiny (headless)
\;$\rightarrow$\; 768-D features \;$\rightarrow$\; DANN/GRL \;$\rightarrow$\;
Autoencoder ($768\rightarrow64$) \;$\rightarrow$\;
6-qubit VQC (54 params) \;$\rightarrow$\; $P(\text{pneumonia})$
\vspace{1.2em}}}
\caption{Schematic of the hybrid CNN--autoencoder--VQC pipeline.}
\label{fig:pipeline}
\end{figure}

\subsection{Dataset and Preprocessing}

The pipeline was evaluated on the publicly available pediatric CXR pneumonia
dataset originating from the Guangzhou Women and Children's Medical Center
(5,856 anteroposterior CXR images, ages approximately 1--5 years).  The
original release provides a validation split of only 16 images, which is
insufficient for reliable convergence monitoring; the training and validation
partitions were therefore merged and re-split 80:20 with class stratification,
while the original test partition was left untouched to preserve the integrity
of the final evaluation.  The resulting split is summarized in
Table~\ref{tab:dataset}.  Notably, the test set exhibits a materially different
class prevalence (62.5\% pneumonia) than the training and validation sets
(74.2\% pneumonia), constituting an explicit, quantified dataset shift that is
exploited below as a robustness stress test rather than concealed.

\begin{table}[htbp]
\centering
\caption{Dataset composition and class-prevalence shift.}
\label{tab:dataset}
\begin{tabular}{@{}lrrr@{}}
\toprule
Split & Images & Pneumonia (\%) & Normal (\%) \\
\midrule
Training   & 4{,}185 & 74.2 & 25.8 \\
Validation & 1{,}047 & 74.2 & 25.8 \\
Test       &   624   & 62.5 & 37.5 \\
\bottomrule
\end{tabular}
\end{table}

Standard image preprocessing was applied prior to the ConvNeXt-Tiny forward
pass: resizing to $224\times224$ pixels, conversion to a normalized float
tensor, and channel-wise standardization using ImageNet mean $[0.485, 0.456,
0.406]$ and standard deviation $[0.229, 0.224, 0.225]$, matching the
statistics of the pretraining distribution.  Because the CXR images are
grayscale while the backbone expects three input channels, the single
channel is duplicated across the RGB channels.  Medical-safe augmentations
were limited to mild rotations ($\pm 7^\circ$), small affine translations,
and modest color jitter; aggressive augmentations such as horizontal flipping
and RandAugment were deliberately excluded as they are harmful for
anatomical X-ray content.

\subsection{Training Configuration and Loss Functions}

The VQC and a classical baseline classifier (a single-hidden-layer MLP with
32 hidden units, operating on the same 64-dimensional autoencoder features)
were trained independently, each with the Adam optimizer, a batch size of 16,
and early stopping (patience of 3 epochs) on validation loss.  Because VQC
outputs (Equation~\ref{eq:vqc}) lie in $[-1,1]$ rather
than $[0,1]$, binary labels $y_i \in \{0,1\}$ were remapped to targets
$t_i \in \{-1,+1\}$, and the quantum classifier was trained with a
class-weighted mean squared error,
\begin{equation}
  \mathcal{L}_{\text{QMSE}} = \frac{1}{N}\sum_{i=1}^{N}
  w(t_i)\left(z_i - t_i\right)^2,
  \label{eq:qmse}
\end{equation}
where $z_i$ is the raw circuit output and $w(\cdot)$ up-weights the minority
(Normal) class to counteract the training-set prevalence imbalance.  The
classical baseline was trained with a class-weighted binary cross-entropy loss
on its sigmoid output.  The quantum classifier's learning rate followed a
linear warm-up over the first $E_{\text{warmup}} = 3$ epochs followed by
cosine annealing,
\begin{equation}
  \eta(e) =
  \begin{cases}
    \eta_0 \dfrac{e}{E_{\text{warmup}}}, & e \le E_{\text{warmup}}, \\[8pt]
    \eta_{\min} + \dfrac{1}{2}(\eta_0 - \eta_{\min})\!\left[1 +
    \cos\!\left(\pi\, \dfrac{e - E_{\text{warmup}}}{E - E_{\text{warmup}}}
    \right)\right], & e > E_{\text{warmup}},
  \end{cases}
  \label{eq:cosine_schedule}
\end{equation}
with $\eta_0 = 10^{-3}$ and $\eta_{\min} = 10^{-5}$, while the classical
baseline used unmodified cosine annealing.  Gradient computation for the
quantum circuit used the adjoint differentiation method on the simulator,
which is dramatically faster than the parameter-shift rule and requires no
additional circuit executions beyond the forward pass.  All quantum-circuit
simulation was performed on a GPU-accelerated statevector simulator
(noise-free), so the reported results characterize the idealized,
decoherence-free behavior of the ansatz rather than performance on physical
NISQ hardware.  In practice both models converged well before the 50-epoch
budget (early stopping triggered at approximately 7--10 epochs).

\subsection{Results}

Table~\ref{tab:results} summarizes test-set performance for the classical
baseline and the hybrid VQC classifier at the validation-selected decision
thresholds (Section~\ref{sec:threshold}): $\tau=0.30$ for the MLP and
$\tau=0.35$ for the VQC.  Both thresholds were selected to maximize balanced
accuracy on the validation set only, and applied once to the test set.

\begin{table}[htbp]
\centering
\caption{Test-set diagnostic performance: classical MLP baseline vs.\ hybrid
  VQC classifier (624 test images, 62.5\% pneumonia prevalence).  Thresholds
  $\tau$ were selected on the validation set via balanced-accuracy
  maximization and applied once to the test set.}
\label{tab:results}
\begin{tabular}{@{}lrr@{}}
\toprule
Metric & Classical (MLP) & Hybrid (VQC) \\
\midrule
Accuracy              & 82.53\% & 81.25\% \\
Precision             & 78.27\% & 77.03\% \\
Sensitivity (Recall)  & 99.74\% & 99.74\% \\
Specificity           & 53.85\% & 50.43\% \\
$F_1$-score           & 0.8771  & 0.8693 \\
AUC-ROC (test)        & 0.940   & 0.860  \\
Balanced accuracy     & 0.7679  & 0.7509 \\
\addlinespace
Decision threshold $\tau$ & 0.30 & 0.35 \\
Decision-layer parameters & 2{,}113 & 54 \\
Parameter reduction       & \multicolumn{2}{c}{$\approx 39\times$} \\
\bottomrule
\end{tabular}
\end{table}

The classical baseline outperforms the hybrid quantum classifier on every
reported metric, consistent with the expectation that a fully connected head
can better exploit a rich, 64-dimensional continuous representation than a
shallow, locally entangling 6-qubit circuit.  The gap is nevertheless
modest---the accuracy difference is 1.28 percentage points and is not
statistically significant at $\alpha=0.05$ under McNemar's test applied to
the 624 paired test predictions---so the hybrid model attains essentially
comparable discriminative performance (test AUC-ROC 0.860) using
approximately $39\times$ fewer decision-layer parameters (54 versus 2{,}113).
Both models exhibit a ``cautious screening'' operating profile, with very
high sensitivity (99.74\%) at the cost of low specificity (50.4--53.9\%),
which from a purely metric standpoint favors minimizing missed pneumonia
cases over minimizing false alarms; we emphasize that no clinical validation
of this trade-off is implied.  The modest specificity is attributable to the
decision thresholds selected by the balanced-accuracy criterion, which
deliberately trades specificity for sensitivity in an imbalanced test
population.

Compared to the literature surveyed in Table~\ref{tab:related}, the test
AUC-ROC of 0.860 is lower than the near-perfect figures reported by Rao et
al.~\cite{rao2024} or Kulkarni et al.~\cite{kulkarni2023}, differences
that can be attributed to: (i)~the explicit train/test class-prevalence
shift applied here; (ii)~the single-run evaluation without cross-validation;
and (iii)~the absence of fine-tuning of the ConvNeXt-Tiny backbone.  These
factors collectively reduce the upper bound on achievable performance and
make the comparison more conservative and, arguably, more realistic.

% ==============================================================
\section{Experiments on Real Quantum Hardware}
\label{sec:hardware}
% ==============================================================

\hwplaceholder{%
  This section will report the results of executing the VQC classifier
  (Section~\ref{sec:hybrid}) on one or more physical NISQ devices (e.g.,
  IBM~Quantum superconducting processors or trapped-ion platforms).  It will
  include the following subsections once the experiments are complete.
}

\subsection{Motivation and Hardware Selection}
\label{sec:hw_motivation}

\hwplaceholder{%
  Justification for the choice of hardware platform(s); qubit counts, native
  gate sets, and typical two-qubit gate error rates of the selected devices;
  circuit transpilation and optimization strategy (e.g., routing, SWAP
  insertion, basis-gate decomposition); and the mapping between the 6-qubit
  VQC ansatz and the device's connectivity graph.
}

\subsection{Error Mitigation Strategies}
\label{sec:hw_mitigation}

\hwplaceholder{%
  Description of the error mitigation protocol(s) to be applied:
  readout-error calibration, zero-noise extrapolation (ZNE), probabilistic
  error cancellation (PEC), or dynamical decoupling (DD).  Discussion of
  encoding-dependent mitigation effectiveness in light of the comparative
  literature \cite{qkernel2022,haug2021}, and the expected trade-off between
  mitigation overhead (additional circuit executions) and performance gain.
}

\subsection{Hardware Results and Simulator--Hardware Gap}
\label{sec:hw_results}

\hwplaceholder{%
  Test-set performance metrics (Table~\ref{tab:results} format) obtained on
  the physical device, with and without error mitigation.  Quantitative
  comparison against the noise-free statevector simulator results to assess
  the magnitude of the hardware degradation.  Analysis of which circuit
  components (state preparation, entangling layers, measurement) contribute
  most to the performance gap, and recommendations for circuit re-design to
  improve hardware fidelity.
}

% ==============================================================
\section{Experimental Considerations and Evaluation Metrics}
\label{sec:eval}
% ==============================================================

\subsection{Metrics}

Given the class imbalance typical of public pneumonia CXR datasets, accuracy
alone is an insufficient evaluation criterion.  We recommend reporting, at
minimum: sensitivity (recall), $\mathrm{Sens} = TP/(TP+FN)$; specificity,
$\mathrm{Spec} = TN/(TN+FP)$; balanced accuracy,
$\mathrm{BalAcc} = \tfrac{1}{2}(\mathrm{Sens}+\mathrm{Spec})$; the
$F_1$-score; and the area under the receiver operating characteristic curve
(AUC-ROC), which is threshold-independent and therefore particularly
informative when comparing classifiers (such as a VQC and an MLP) whose raw
outputs have different natural ranges and calibration properties.

\subsection{Threshold Selection}
\label{sec:threshold}

Because the VQC output (Equation~\ref{eq:sigmoid_remap}) and the classical
sigmoid output are continuous scores, a decision threshold $\tau$ must be
chosen to obtain a binary prediction
$\hat{y} = \mathbb{1}[\,p_{\text{pneumonia}} > \tau\,]$.  A principled
approach is to perform a discrete threshold scan on the validation set,
selecting $\tau^\star = \arg\max_\tau \mathrm{BalAcc}(\tau)$, and to report
test-set metrics at this fixed, validation-selected threshold rather than at
an arbitrarily chosen value.  In the case study reviewed here, this procedure
selected $\tau^\star = 0.30$ for the MLP and $\tau^\star = 0.35$ for the VQC
over a scan range of $\tau \in [0.30, 0.80]$ with step $0.025$, both below
the conventional default of 0.5.  The sensitivity analysis itself remains a
useful diagnostic: thresholds near the conventional 0.5 default shift the
operating point toward higher specificity at the cost of sensitivity, whereas
the balanced-accuracy criterion favors the low-threshold, high-sensitivity
operating point reported in Table~\ref{tab:results}.

\subsection{Dataset Shift as an Explicit Evaluation Axis}

We highlight dataset shift as a first-class evaluation concern rather than
an incidental property of the data split.  Reporting class prevalence
separately for training, validation, and test partitions (as in
Table~\ref{tab:dataset}) allows a reader to distinguish genuine
generalization performance from performance inflated by matching train/test
class ratios.  This is particularly relevant for QML studies, where
sample-efficiency constraints (imposed by simulator run-time or hardware queue
availability) often preclude the large-scale cross-validation protocols that
are standard in classical deep learning benchmarking; a single, explicitly
characterized shift-aware split is a partial, low-cost mitigation.  The
generalizability limitations of single-institution, narrow-demographic datasets
highlighted by Zech et al.~\cite{zech2018} apply with equal force to hybrid
quantum--classical architectures.

\subsection{Interpretability}

For classical CNN components, gradient-based saliency methods such as Grad-CAM
can verify that the network attends to clinically relevant lung regions rather
than to background artifacts (e.g., imaging equipment, laterality markers),
providing a qualitative sanity check complementary to aggregate performance
metrics.  Decoodt et al.~\cite{decoodt2023} found that augmenting a
DenseNet-121 backbone with a variational quantum layer improved
Grad-CAM$++$ localization quality for cardiomegaly detection, suggesting a
potential ancillary benefit of quantum components beyond raw classification
accuracy.  Analogous interpretability tools for the quantum component are
considerably less mature; feasible directions include measuring the sensitivity
of the circuit output to perturbations of individual input amplitudes, or
analyzing feature importance in the classical dimensionality-reduction stage
(e.g., via SelectKBest or PCA loading vectors) as a proxy for what information
reaches the quantum classifier.

% ==============================================================
\section{Challenges, Limitations, and Open Problems}
\label{sec:limitations}
% ==============================================================

We summarize the principal limitations of current hybrid VQC approaches to CXR
pneumonia detection, several of which are directly evidenced by the case study
in Section~\ref{sec:hybrid}.

\begin{enumerate}[leftmargin=1.4em]
  \item \textbf{Barren plateaus.}  As circuit width or entangling depth
    increases, gradient variance decays exponentially
    (Equation~\ref{eq:barren}), limiting the ansätze that can be reliably
    trained with current optimizers \cite{mcclean2018,cerezo2021,marrero2021}.
    The case-study ansatz was therefore deliberately restricted to three
    shallow layers on six qubits, a regime in which gradient-variance
    concentration is not expected to dominate training
    (Section~\ref{sec:barren}).  While several mitigation strategies exist,
    none fully eliminates the fundamental exponential scaling
    \cite{kashif2023}.
  \item \textbf{State-preparation and encoding overhead.}  Amplitude encoding
    achieves exponential qubit compression but, in the general case, incurs
    exponential state-preparation circuit depth \cite{aaronson2015};
    simulator-based studies typically do not account for this cost, which will
    materially affect feasibility on physical hardware.
  \item \textbf{Dataset shift.}  Differences in class prevalence (and, more
    broadly, in acquisition device, patient demographics, and imaging protocol)
    between training and deployment distributions are a well-documented source
    of generalization failure for CXR classifiers \cite{zech2018}, and were
    observed to degrade specificity in the reviewed case study.
  \item \textbf{Hardware noise.}  All simulation results reviewed here were
    obtained on noise-free statevector simulators; real NISQ devices introduce
    gate errors, readout errors, and decoherence that are expected to degrade
    performance unless mitigated.  Empirical evidence from quantum kernel
    studies suggests that error mitigation effectiveness is strongly
    encoding-dependent and that no universally applicable mitigation protocol
    has emerged \cite{qkernel2022,haug2021}.
  \item \textbf{Information loss under dimensionality reduction.}  Compressing
    a 768-dimensional CNN embedding to 64 (or fewer) dimensions via the
    nonlinear autoencoder
    necessarily discards variance, some fraction of which may be diagnostically
    relevant.  Alternative encodings that avoid aggressive compression
    (e.g., data re-uploading \cite{perez2020} with angle encoding, or
    distributed amplitude encoding over multiple shorter circuits) remain open
    design questions \cite{sharma2023survey}.
  \item \textbf{Absence of demonstrated quantum advantage.}  At the circuit
    scales considered (single-digit qubits), no computational or statistical
    advantage over classical models has been established; the case for VQC
    decision layers currently rests on parameter efficiency rather than
    accuracy or speed.
  \item \textbf{Limited statistical robustness.}  The case study reviewed here
    uses a single train/validation/test split without cross-validation or
    repeated randomized resampling, owing to the computational cost of circuit
    simulation;
    reported performance differences, while accompanied by approximate
    confidence intervals, should be interpreted with the caveats appropriate
    to a single-run evaluation.
  \item \textbf{Single-institution, narrow-demographic data.}  The underlying
    dataset originates from a single pediatric population at one medical center,
    limiting claims of generalizability to adult populations, other acquisition
    protocols, or other institutions.
  \item \textbf{Unfair classical baselines.}  Several studies surveyed in
    Table~\ref{tab:related} compare a quantum model favorably against a
    deliberately weak classical reference, inflating the apparent advantage of
    the hybrid approach.  Rigorous evaluation requires comparison against
    optimized classical baselines of equal or greater parameter budget on the
    quantum side \cite{aaronson2015}.
\end{enumerate}

% ==============================================================
\section{Future Directions}
\label{sec:future}
% ==============================================================

Several directions appear promising for advancing hybrid quantum--classical
pneumonia detection beyond the proof-of-concept stage.

\begin{itemize}[leftmargin=1.4em]
  \item \textbf{Execution on physical quantum hardware.}  Migrating the VQC
    component to a cloud-accessible NISQ device (e.g., superconducting or
    trapped-ion platforms), combined with hardware-aware error mitigation
    \cite{qkernel2022,haug2021}, will quantify the simulator-to-hardware
    performance gap and is a necessary step before any hardware-relevant
    advantage claim can be assessed.  This is the primary goal of
    Section~\ref{sec:hardware} (forthcoming).
  \item \textbf{Quantum convolutional and quanvolutional architectures.}
    QCNNs \cite{cong2019,wei2022} that exploit spatial locality in the quantum
    circuit, and quanvolutional pre-processing layers \cite{henderson2020quanvolutional}
    that apply quantum filters to image patches prior to classical CNN
    processing, represent theoretically motivated alternatives to the
    ``quantum classifier head'' paradigm.  Their trainability guarantees
    \cite{pesah2021} and parameter independence from input size make them
    particularly attractive for deployment under NISQ constraints.
  \item \textbf{Quantum kernel methods.}  Replacing the VQC classifier with a
    quantum kernel SVM \cite{havlicek2019,schuld2019kernel} has the advantage
    of convex optimization---eliminating barren plateaus by construction---and
    a well-defined notion of kernel quality via the quantum Fisher information
    \cite{haug2021}.  Scalable kernel estimation via randomized measurements
    \cite{haug2021} makes this approach computationally feasible even for
    training sets of several thousand images.
  \item \textbf{Data re-uploading classifiers.}  Applying the
    data re-uploading paradigm \cite{perez2020} within the VQC decision layer
    would increase circuit expressibility without requiring additional qubits,
    potentially alleviating the information-loss constraints imposed by
    aggressive dimensionality reduction to the 64-dimensional amplitude-encoded
    register.
  \item \textbf{Systematic hyperparameter optimization.}  Bayesian optimization
    frameworks (e.g., Optuna) applied jointly across classical and quantum
    hyperparameters (learning rate, ansatz depth, entangling topology, number
    of retained latent dimensions) would allow more efficient exploration of the
    design space than manual or grid search, particularly given the high per-run
    cost of circuit simulation.
  \item \textbf{Rigorous, shift-aware, multi-institutional evaluation.}
    Extending evaluation to larger, multi-site datasets (e.g., CheXpert,
    NIH ChestX-ray14) and explicit domain-adaptation techniques (reweighting,
    adversarial domain adaptation) would provide a more realistic assessment
    of generalization than single-institution pediatric data \cite{zech2018}.
  \item \textbf{Ensemble and hybrid-decision strategies.}  Combining classical
    and quantum model outputs (e.g., via weighted logit averaging or a
    meta-classifier) could exploit differences in the error profiles of the two
    approaches, potentially improving both sensitivity and specificity relative
    to either model alone \cite{kundu2021}.
  \item \textbf{Quantum-side interpretability.}  Developing
    sensitivity- and perturbation-based tools analogous to Grad-CAM for the
    quantum decision layer would help clarify whether the VQC learns
    complementary or redundant information relative to the classical
    autoencoder projection, and would support the broader goal of building
    trustworthy clinical decision-support systems.
\end{itemize}

% ==============================================================
\section{Conclusion}
\label{sec:conclusion}
% ==============================================================

This paper reviewed the computational and algorithmic foundations of hybrid
quantum--classical architectures for pneumonia detection from chest X-ray
images, with an emphasis on data-encoding strategies, variational and
convolutional ansatz design, quantum kernel methods, and the trainability
limitations imposed by barren plateaus together with their principal
mitigation strategies.  We surveyed a growing body of literature applying
hybrid QML pipelines to CXR tasks---including recent works demonstrating
statistically significant improvements over classical baselines
\cite{kulkarni2023} and near-unity accuracy on COVID-19 radiography datasets
\cite{rao2024}---while noting the evaluation limitations that temper direct
comparison across studies.

Using a ConvNeXt-Tiny/autoencoder/VQC pipeline evaluated on a public pediatric
CXR dataset as a concrete case study, we showed that a 6-qubit, 54-parameter
variational quantum classifier can achieve a level of diagnostic performance
comparable to---though not exceeding---that of a classical MLP of comparable
input dimensionality (test AUC-ROC 0.860 versus 0.940; accuracy 81.25\%
versus 82.53\%), while requiring approximately $39\times$ fewer trainable
parameters in the decision layer (54 versus 2{,}113).  This result
should be read narrowly: it demonstrates that a compact variational circuit can
serve as a functional, parameter-efficient classifier on a compressed
medical-imaging representation, not that a computational or statistical quantum
advantage has been achieved.  Barren plateaus, state-preparation overhead,
aggressive dimensionality-reduction-induced information loss, and dataset shift
all constrain current performance, and the simulation-based results reported
here have yet to be validated on physical quantum hardware.

We regard hybrid VQC decision layers, quantum convolutional architectures, and
quantum kernel classifiers as promising but early-stage research directions
whose eventual clinical or practical relevance will depend on: progress in
fault-tolerant or higher-fidelity NISQ hardware; more expressive yet trainable
ansatz designs (e.g., data re-uploading, low-entanglement initialization);
hardware-aware error mitigation; and substantially more rigorous,
multi-institutional, shift-aware evaluation protocols than are currently
standard in the literature.  The forthcoming hardware experiments
(Section~\ref{sec:hardware}) will provide the first direct empirical
quantification of the simulator-to-hardware gap for this specific pipeline and
will constitute a necessary precondition for any hardware-grounded advantage
claim.

% ==============================================================
%  REFERENCES
% ==============================================================
\begin{thebibliography}{99}

\bibitem{aaronson2015}
Aaronson, S. (2015). Read the fine print. \emph{Nature Physics}, 11(4),
291--293.

\bibitem{biamonte2017}
Biamonte, J., Wittek, P., Pancotti, N., Rebentrost, P., Wiebe, N., \& Lloyd,
S. (2017). Quantum machine learning. \emph{Nature}, 549(7671), 195--202.

\bibitem{cdc2025}
Centers for Disease Control and Prevention. (2025). \emph{About Pneumonia}.
\url{https://www.cdc.gov/pneumonia/about/index.html}

\bibitem{cerezo2021}
Cerezo, M., Sone, A., Volkoff, T., Cincio, L., \& Coles, P.~J. (2021).
Cost-function-dependent barren plateaus in shallow parametrized quantum
circuits. \emph{Nature Communications}, 12(1), 1791.

\bibitem{cong2019}
Cong, I., Choi, S., \& Lukin, M.~D. (2019). Quantum convolutional neural
networks. \emph{Nature Physics}, 15(12), 1273--1278.

\bibitem{decoodt2023}
Decoodt, P., Liang, T.~J., et al. (2023). Hybrid classical--quantum transfer
learning for cardiomegaly detection in chest X-rays. \emph{Journal of
Imaging}, 9(7), 128.
\url{https://doi.org/10.3390/jimaging9070128}

\bibitem{ding2024}
Ding, Y., Zhang, S., \& Li, X. (2024). Noise mitigation: an efficient approach
for quantum convolutional neural networks to enhance performance in NISQ.
Preprint. \url{https://doi.org/10.21203/rs.3.rs-4594670/v1}

\bibitem{haug2021}
Haug, T., Self, C.~N., \& Kim, M. (2021). Quantum machine learning of
large datasets using randomized measurements. \emph{arXiv preprint
arXiv:2108.01039}.

\bibitem{havlicek2019}
Havl\'i\v{c}ek, V., C\'orcoles, A.~D., Temme, K., Harrow, A.~W., Kandala,
A., Chow, J.~M., \& Gambetta, J.~M. (2019). Supervised learning with
quantum-enhanced feature spaces. \emph{Nature}, 567(7747), 209--212.

\bibitem{he2016}
He, K., Zhang, X., Ren, S., \& Sun, J. (2016). Deep residual learning for
image recognition. In \emph{Proceedings of the IEEE Conference on Computer
Vision and Pattern Recognition (CVPR)}, 770--778.

\bibitem{henderson2020quanvolutional}
Henderson, M., Shakya, S., Pradhan, S., \& Cook, T. (2020). Quanvolutional
neural networks: Powering image recognition with quantum circuits.
\emph{Quantum Machine Intelligence}, 2(1), 2.

\bibitem{iqwig2021}
Institute for Quality and Efficiency in Health Care (IQWiG). (2021).
\emph{Overview: Pneumonia}. InformedHealth.org.
\url{https://www.ncbi.nlm.nih.gov/books/NBK525774/}

\bibitem{kashif2023}
Kashif, M., \& Al-Kuwari, S. (2023). The unified effect of data encoding,
ansatz expressibility and entanglement on the trainability of HQNNs.
\emph{International Journal of Parallel, Emergent and Distributed Systems},
38(5), 362--400.
\url{https://doi.org/10.1080/17445760.2023.2231163}

\bibitem{kulkarni2023}
Kulkarni, V., Pawale, S., \& Kharat, A. (2023). A classical--quantum
convolutional neural network for detecting pneumonia from chest radiographs.
\emph{Neural Computing and Applications}, 35(21), 15503--15510.
\url{https://doi.org/10.1007/s00521-023-08566-1}

\bibitem{khatoniar2024}
Khatoniar, R., Vyas, V., Acharya, V., et al. (2024). Quantum convolutional
neural network for medical image classification. In \emph{Proceedings of the
2nd International Conference on Trends in Quantum Computing and Emerging
Business Technologies (TQCEBT)}. IEEE.
\url{https://doi.org/10.1109/tqcebt59414.2024.10545121}

\bibitem{kundu2021}
Kundu, R., Das, R., Geem, Z.~W., Han, G.-T., \& Sarkar, R. (2021). Pneumonia
detection in chest X-ray images using an ensemble of deep learning models.
\emph{PLOS ONE}, 16(9), e0256630.

\bibitem{liu2022sea}
Liu, X., Liu, G., Huang, J., et al. (2022). Mitigating barren plateaus of
variational quantum eigensolvers. \emph{arXiv preprint arXiv:2205.13539}.

\bibitem{marrero2021}
Ortiz Marrero, C., Kieferov\'a, M., \& Wiebe, N. (2021).
Entanglement-induced barren plateaus. \emph{PRX Quantum}, 2(4), 040316.
\url{https://doi.org/10.1103/PRXQuantum.2.040316}

\bibitem{mcclean2018}
McClean, J.~R., Boixo, S., Smelyanskiy, V.~N., Babbush, R., \& Neven, H.
(2018). Barren plateaus in quantum neural network training landscapes.
\emph{Nature Communications}, 9(1), 4812.

\bibitem{mitarai2018}
Mitarai, K., Negoro, M., Kitagawa, M., \& Fujii, K. (2018). Quantum circuit
learning. \emph{Physical Review A}, 98(3), 032309.

\bibitem{naz2020}
Naz, A. (2020). Literature review on X-ray based pneumonia detection using
deep learning. \emph{LCJSTEM}, 1(1), 1--9.

\bibitem{nielsen2010}
Nielsen, M.~A., \& Chuang, I.~L. (2010). \emph{Quantum Computation and
Quantum Information} (10th anniversary ed.). Cambridge University Press.

\bibitem{patti2021}
Patti, T.~L., Najafi, K., Gao, X., \& Yelin, S.~F. (2021). Entanglement
devised barren plateau mitigation. \emph{Physical Review Research}, 3,
033090. \url{https://doi.org/10.1103/PhysRevResearch.3.033090}

\bibitem{perez2020}
P\'erez-Salinas, A., Cervera-Lierta, A., Gil-Fuster, E., \& Latorre, J.~I.
(2020). Data re-uploading for a universal quantum classifier. \emph{Quantum},
4, 226.

\bibitem{pesah2021}
Pesah, A., Cerezo, M., Wang, S., Volkoff, T., Sornborger, A.~T., \& Coles,
P.~J. (2021). Absence of barren plateaus in quantum convolutional neural
networks. \emph{Physical Review X}, 11(4), 041011.

\bibitem{preskill2018}
Preskill, J. (2018). Quantum computing in the NISQ era and beyond.
\emph{Quantum}, 2, 79.

\bibitem{qkernel2022}
Beaulieu, D., Miracle, D., Pham, A.~T., \& Scherr, W. (2022). Quantum kernel
for image classification of real-world manufacturing defects. \emph{arXiv
preprint arXiv:2212.08693}.

\bibitem{rao2024}
Rao, G.~V.~E., Rajitha, B., Srinivasu, P.~N., Ijaz, M.~F., \& Wo\'zniak, M.
(2024). Hybrid framework for respiratory lung diseases detection based on
classical CNN and quantum classifiers from chest X-rays. \emph{Biomedical
Signal Processing and Control}, 88, 105567.
\url{https://doi.org/10.1016/j.bspc.2023.105567}

\bibitem{schuld2018}
Schuld, M., \& Petruccione, F. (2018). \emph{Supervised Learning with
Quantum Computers}. Springer, Quantum Science and Technology.

\bibitem{schuld2019kernel}
Schuld, M., \& Killoran, N. (2019). Quantum machine learning in feature
Hilbert spaces. \emph{Physical Review Letters}, 122(4), 040504.

\bibitem{schuld2021kernel}
Schuld, M. (2021). Supervised quantum machine learning models are kernel
methods. \emph{arXiv preprint arXiv:2101.11020}.

\bibitem{ortiz2022}
Ortiz-Toro, C., Garc\'ia-Pedrero, A., Lillo-Saavedra, M., \& Gonzalo-Mart\'in,
C. (2022). Automatic detection of pneumonia in chest X-ray images using
textural features. \emph{Computers in Biology and Medicine}, 145, 105466.
\url{https://doi.org/10.1016/j.compbiomed.2022.105466}

\bibitem{sharma2023survey}
Sharma, A., et al. (2023). Survey of encoding techniques for quantum machine
learning. \emph{Cybernetics and Physics}, 12(1), 45--67.

\bibitem{sharmaR2023hybrid}
Sharma, R. (2023). Classification of pneumonia via a hybrid ZFNet--quantum
neural network. \emph{Current Medical Imaging Reviews}. Advance online
publication.

\bibitem{sim2019}
Sim, S., Johnson, P.~D., \& Aspuru-Guzik, A. (2019). Expressibility and
entangling capability of parametrized quantum circuits for hybrid
quantum-classical algorithms. \emph{Advanced Quantum Technologies}, 2(12),
1900070.

\bibitem{wei2022}
Wei, S., Chen, Y., Zhou, Z.-R., \& Long, G.-L. (2022). A quantum convolutional
neural network on NISQ devices. \emph{AAPPS Bulletin}, 32(1), 2.
\url{https://doi.org/10.1007/s43673-021-00030-3}

\bibitem{zech2018}
Zech, J.~R., Badgeley, M.~A., Liu, M., Costa, A.~B., Titano, J.~J., \&
Oermann, E.~K. (2018). Variable generalization performance of a deep learning
model to detect pneumonia in chest radiographs: A cross-sectional study.
\emph{PLOS Medicine}, 15(11), e1002683.

\bibitem{zhu2024}
Zhu, Y., Liang, Y., \& Situ, H. (2024). Enhancing variational quantum circuit
training: an improved neural network approach for barren plateau mitigation.
\emph{arXiv preprint arXiv:2411.09226}.

\end{thebibliography}

\end{document}
